\documentclass[conference]{IEEEtran}
\IEEEoverridecommandlockouts

% Configuración para español
\usepackage[spanish,es-nodecimaldot,es-noshorthands]{babel}
\usepackage[utf8]{inputenc}
\usepackage[T1]{fontenc}

% Paquetes estándar IEEE
\usepackage{cite}
\usepackage{amsmath,amssymb,amsfonts}
\usepackage{graphicx}
\graphicspath{{./}{Evidencias/}{img/}{../avance3/img/}}
\usepackage{textcomp}
\usepackage{xcolor}
\usepackage{float}
\usepackage{listings}
\usepackage{multirow}
\usepackage{array}
\usepackage[hidelinks]{hyperref}

% Colores y macros de apoyo
\definecolor{codegray}{rgb}{0.45,0.45,0.45}
\definecolor{codeblue}{rgb}{0.10,0.10,0.60}
\definecolor{pend}{rgb}{0.75,0.10,0.10}

\lstset{
	basicstyle=\ttfamily\scriptsize,
	keywordstyle=\color{codeblue}\bfseries,
	commentstyle=\color{codegray}\itshape,
	stringstyle=\color{codegray},
	frame=single,
	breaklines=true,
	breakatwhitespace=true,
	showstringspaces=false,
	columns=fullflexible,
	tabsize=2,
	literate=%
	{á}{{\'a}}1 {é}{{\'e}}1 {í}{{\'i}}1 {ó}{{\'o}}1 {ú}{{\'u}}1
	{Á}{{\'A}}1 {É}{{\'E}}1 {Í}{{\'I}}1 {Ó}{{\'O}}1 {Ú}{{\'U}}1
	{ñ}{{\~n}}1 {Ñ}{{\~N}}1 {¿}{{?`}}1 {¡}{{!`}}1
}

\newcommand{\imgph}[2]{%
	\fbox{\begin{minipage}[c][#2][c]{0.92\columnwidth}
			\centering\textcolor{codegray}{\footnotesize [Espacio para imagen: #1]}
	\end{minipage}}%
}
% Marca en rojo los valores medidos pendientes; eliminar antes de la entrega final
\newcommand{\pend}[1]{\textcolor{pend}{[#1]}}

\def\BibTeX{{\rm B\kern-.05em{\sc i\kern-.025em b}\kern-.08em
	T\kern-.1667em\lower.7ex\hbox{E}\kern-.125emX}}

\begin{document}

\title{Reporte Final de Proyecto: \\
	\large Terminal de Juegos Retro Distribuida con Mando Embebido}

\author{
	\IEEEauthorblockN{Belmont Muñoz Samuel}
	\IEEEauthorblockA{\textit{Ingeniería en Computación} \\
		\textit{Facultad de Ingeniería, UNAM}\\
		No. Cuenta: 317318289}
	\and
	\IEEEauthorblockN{Domínguez Palacios Jesús Alejandro}
	\IEEEauthorblockA{\textit{Ingeniería en Computación} \\
		\textit{Facultad de Ingeniería, UNAM}\\
		No. Cuenta: 421066205}
	\and
	\IEEEauthorblockN{Dorantes Mejía Emmanuel}
	\IEEEauthorblockA{\textit{Ingeniería en Computación} \\
		\textit{Facultad de Ingeniería, UNAM}\\
		No. Cuenta: 313004766}
	\and
	\IEEEauthorblockN{García Olvera Emiliano}
	\IEEEauthorblockA{\textit{Ingeniería en Computación} \\
		\textit{Facultad de Ingeniería, UNAM}\\
		No. Cuenta: 314256009}
	\and
	\IEEEauthorblockN{Hernández Alonso Antonio Adonahí}
	\IEEEauthorblockA{\textit{Ingeniería en Computación} \\
		\textit{Facultad de Ingeniería, UNAM}\\
		No. Cuenta: 316339825}
	\and
	\IEEEauthorblockN{Jiménez Reyes Alan Yael}
	\IEEEauthorblockA{\textit{Ingeniería en Computación} \\
		\textit{Facultad de Ingeniería, UNAM}\\
		No. Cuenta: 320026726}
}

\maketitle

\begin{center}
	\vspace{-0.3cm}
	{\small
		\textbf{Fundamentos de Sistemas Embebidos — Grupo 6, Equipo 2} \quad | \quad
		\textbf{Entrega Final:} Noviembre 2026
	}
	\vspace{0.2cm}
\end{center}

\begin{abstract}
Este documento presenta el diseño, implementación y validación experimental de una terminal de juegos retro distribuida basada en hardware embebido. La arquitectura divide responsabilidades entre un nodo esclavo determinista (Raspberry Pi Pico, RP2040) ejecutando firmware \textit{bare-metal} y un nodo maestro (Raspberry Pi 4 Model B) corriendo RetroArch sobre Raspberry Pi OS Lite. El mando embebido adquiere 10 botones digitales mediante dos registros de desplazamiento 74HC165 en cascada con filtrado RC por hardware, gestionados por una máquina de estados PIO y canales DMA encadenados a $292.6\text{ Hz}$ ($3.417\text{ ms}$), junto con ejes analógicos mediante ADC a $4\text{ kHz}$ por canal. La comunicación con la consola opera mediante un enlace único USB HID bidireccional: reportes de entrada de 10 bytes cada $1\text{ ms}$ y reportes de salida de 8 bytes para gobernar actuadores hápticos, sonoros y visuales sin bloquear el lazo principal. En la consola, el despliegue gráfico se realiza por SPI a $50\text{ MHz}$ sobre una pantalla TFT de 3.5 pulgadas mediante DRM/KMS ($18.3\text{ FPS}$ medidos), mientras que un esquema de comunicación entre procesos (UDP local y socket \texttt{AF\_UNIX}) identifica el título activo, extrae telemetría de memoria en tiempo real y controla los actuadores bajo orquestación autónoma con \texttt{systemd}. Se reportan las métricas experimentales de latencia, uso de CPU, huella de memoria, consumo energético y tolerancia a fallas.
\end{abstract}

\begin{IEEEkeywords}
Sistemas embebidos, Raspberry Pi Pico, Raspberry Pi 4, PIO, DMA, USB HID, RetroArch, systemd, IPC.
\end{IEEEkeywords}

% =====================================================================
% I. INTRODUCCIÓN, PROBLEMA Y ALCANCE
% =====================================================================
\section{Introducción, Problema y Alcance}

\subsection{Descripción del Problema y Usuario Objetivo}
Los sistemas de emulación arcade de bajo costo suelen depender de periféricos genéricos con tiempos de respuesta no deterministas y carecen de retroalimentación física vinculada al estado interno del videojuego. El objetivo de este proyecto es construir una terminal arcade autónoma dirigida a estudiantes y entusiastas de los sistemas embebidos, separando la adquisición determinista de controles físicos de bajo nivel del procesamiento multimedia y de emulación en un sistema operativo embebido.

\subsection{Alcance Implementado y Delimitación}
\begin{itemize}
	\item \textbf{Implementado:} Control físico basado en RP2040 con 10 botones digitales (74HC165 $\times 2$ con antirrebote RC), entrada analógica ADC, actuadores (motor de vibración, buzzer pasivo a $2{,}700\text{ Hz}$ y LEDs), comunicación bidireccional USB HID, emulación en Raspberry Pi 4 sobre pantalla TFT SPI de 3.5 pulgadas, telemetría de memoria multijuego en tiempo real con IPC por socket UNIX y arranque autónomo tolerante a fallas mediante \texttt{systemd}.
	\item \textbf{Ajuste de interfaz de video:} Respecto a la salida externa HDMI a $60\text{ FPS}$ contemplada en el Avance 1, en el prototipo final se integró una pantalla TFT portátil de 3.5'' sobre el bus SPI0 ($50\text{ MHz}$), manteniendo la lógica de emulación interna a $60\text{ Hz}$ con una tasa real de refresco visual en panel de $18.3\text{ FPS}$ ($20.35\text{ FPS}$ teóricos de bus).
	\item \textbf{Fuera de alcance final:} Desarrollo de motores de juego o ROMs propias, multijugador remoto WAN por VPN (WireGuard) y multijugador local en red LAN (declarado explícitamente en la Sección~\ref{sec:requisitos} al priorizar la integración física del gabinete y la telemetría háptica).
\end{itemize}

% =====================================================================
% II. REQUISITOS Y CUMPLIMIENTO
% =====================================================================
\section{Requisitos y Matriz de Cumplimiento}
\label{sec:requisitos}

La Tabla~\ref{tab:requisitos} contrasta los requisitos funcionales (RF) y no funcionales (RNF) definidos al inicio del proyecto frente a los valores medidos en el prototipo final.

\begin{table*}[!t]
	\centering
	\caption{Matriz de cumplimiento de requisitos funcionales y no funcionales}
	\label{tab:requisitos}
	\footnotesize
	\begin{tabular}{|p{1.0cm}|p{3.4cm}|p{3.3cm}|p{3.8cm}|p{2.2cm}|}
		\hline
		\textbf{ID} & \textbf{Requisito} & \textbf{Criterio Planeado} & \textbf{Resultado Medido / Construido} & \textbf{Estado} \\
		\hline
		RF1 & Adquisición de entradas digitales (74HC165 $\times 2$) & Periodo $\le 5.0\text{ ms}$ ($\ge 200\text{ Hz}$) & Periodo de $3.417\text{ ms}$ ($292.6\text{ Hz}$) calculado a partir del programa PIO ($3{,}417$ ciclos, divisor \texttt{125}), $0$ tramas perdidas & \textbf{Cumplido} \\
		\hline
		RF2 & Lectura analógica de eje (ADC) & Resolución $\ge 10\text{ bits}$, periodo $\le 10\text{ ms}$ & Resolución de $12\text{ bits}$ en \textit{round-robin} a $4\text{ kHz}$/canal ($0.25\text{ ms}$ por muestra) con DMA circular & \textbf{Cumplido} \\
		\hline
		RF3 & Telemetría y retroalimentación bidireccional & Recepción de comando en esclavo $< 20\text{ ms}$ tras evento & Lectura UDP (\texttt{127.0.0.1:55355}) + identificación de ROM en \texttt{historial.lpl} + socket UNIX + HID OUT (8 B) con entrega al Pico de $2.2\text{ ms}$ en promedio ($4.1\text{ ms}$ máx., registros del servicio puente) & \textbf{Cumplido} \\
		\hline
		RNF1 & Latencia de control (\textit{Input Lag}) & Extremo a extremo $\le 15\text{ ms}$ & Latencia física medida botón $\rightarrow$ recepción HID en emulador: $7.65\text{ ms}$ en el peor caso (calculado; mejor caso $2.2\text{ ms}$) & \textbf{Cumplido} \\
		\hline
		RNF2 & Respuesta háptica & Activación del motor $\le 50\text{ ms}$ tras señal & Activación PWM a $20\text{ kHz}$ mediante alarma no bloqueante; evento $\rightarrow$ actuador en $23.1\text{ ms}$ típico y $45.8\text{ ms}$ en el peor caso (calculado) & \textbf{Cumplido} \\
		\hline
		RNF3 & Conectividad multijugador en red local (LAN) & Pérdida $0\%$, \textit{ping} $< 5\text{ ms}$ entre dos nodos & No implementado en el prototipo final por disponibilidad de una sola unidad física armada & \textbf{No cumplido (declarado)} \\
		\hline
	\end{tabular}
\end{table*}

% =====================================================================
% III. ARQUITECTURA DEL SISTEMA Y HARDWARE
% =====================================================================
\section{Arquitectura del Sistema y Plataforma}

% Sección escrita por Jesús (diagrama de bloques final, descripción de bloques y pantalla TFT)
% =====================================================================
% Sección de Domínguez Palacios Jesús Alejandro
% Diagrama de bloques final, descripción de bloques y pantalla TFT
% =====================================================================

\subsection{Diagrama de Bloques del Sistema Final}

\begin{figure*}[!t]
	\centering
	% Reemplazar con la ruta correcta al diagrama final
	% \includegraphics[width=\textwidth]{img/diagrama_bloques_final.png}
	\imgph{Diagrama de bloques del sistema completo}{8cm}
	\caption{Arquitectura completa del sistema: nodo esclavo (Pico) y nodo maestro (Pi 4).}
	\label{fig:diagrama_final}
\end{figure*}

\textbf{Descripción pendiente:} Explicar cada bloque del diagrama, flujo de datos entre componentes, y cómo se conectan los dos nodos.

\subsection{Configuración de la Pantalla TFT}

\textbf{Pantalla utilizada:} TFT SPI de 3.5 pulgadas.

\textbf{Configuración pendiente:} Detalles de la configuración en \texttt{config.txt}, driver \texttt{panel-mipi-dbi}, resolución y FPS alcanzados.


\subsection{Justificación de Plataforma Heterogénea y Materiales}
El sistema combina dos paradigmas complementarios:
\begin{enumerate}
	\item \textbf{Nodo Esclavo (Raspberry Pi Pico — RP2040, \textit{Bare-Metal}):} Procesador dual-core ARM Cortex-M0+ a $133\text{ MHz}$ con $264\text{ KB}$ de SRAM y $2\text{ MB}$ de Flash. Sin sistema operativo ni cambios de contexto, las máquinas de estados PIO y los canales DMA ejecutan la adquisición de sensores de forma autónoma, permitiendo que Core 0 gestione el enlace USB HID cada $1\text{ ms}$ y supervise el \textit{watchdog} de hardware ($250\text{ ms}$).
	\item \textbf{Nodo Maestro (Raspberry Pi 4 Model B Rev 1.5, Linux Embebido):} Quad-core ARM Cortex-A72 (hasta $1.8\text{ GHz}$) con $1\text{ GB}$ de RAM LPDDR4 ejecutando Raspberry Pi OS Lite de 64 bits (Debian 13, kernel 6.18) en modo consola sin servidor gráfico, utilizando DRM/KMS para video y \texttt{systemd} para la orquestación de procesos.
\end{enumerate}

El costo total de materiales del prototipo armado asciende a \textbf{\$1,620 MXN} (Raspberry Pi 4 B, Raspberry Pi Pico, dos registros 74HC165, 10 botones arcade con red antirrebote RC de $330\,\Omega$ y $100\text{ nF}$, potenciómetro de $10\text{ k}\Omega$, motor de vibración con etapa NPN y diodo de rueda libre, buzzer pasivo, LEDs con resistencias de $220\,\Omega$ y pantalla TFT SPI de 3.5'').

% =====================================================================
% IV. IMPLEMENTACIÓN, MÉTRICAS Y ROBUSTEZ
% =====================================================================
\section{Implementación, Métricas y Pruebas}

\subsection{Resumen de Métricas Experimentales del Sistema}
La Tabla~\ref{tab:metricas_finales} consolida las cuatro métricas de desempeño medidas directamente sobre el sistema físico operando en condiciones normales de juego (emulación de \textit{Pac-Man}, \textit{Pong} y \textit{Tetris} con telemetría activa).

\begin{table}[H]
	\centering
	\caption{Métricas experimentales medidas en el sistema final}
	\label{tab:metricas_finales}
	\footnotesize
	\begin{tabular}{|p{2.3cm}|p{2.5cm}|p{2.6cm}|}
		\hline
		\textbf{Métrica} & \textbf{Nodo Esclavo (Pico)} & \textbf{Nodo Maestro (Pi 4)} \\
		\hline
		Latencia / Periodo & Muestreo botones: $3.417\text{ ms}$; reporte USB HID: $1.00\text{ ms}$ & Cuadro SPI TFT: $49.15\text{ ms}$ ($18.3\text{ FPS}$ reales); entrada total: $7.65\text{ ms}$ (peor caso calculado) \\
		\hline
		Uso de CPU & $< 2\%$ en Core 0 (PIO + DMA autónomos; \texttt{\_\_wfi()}) & RetroArch: $20.0\%$ de un núcleo; telemetría Python: $1.0\%$; CPU total: $7.6\%$ ($92.4\%$ idle) \\
		\hline
		Uso de Memoria & Flash: $23{,}476\text{ B}$ ($1.1\%$); SRAM: $3{,}520\text{ B}$ ($1.3\%$) & Reposo: $220\text{ MiB}$ / $905\text{ MiB}$ ($24.3\%$); en juego: $283\text{ MiB}$ ($31.3\%$), sin \textit{swap} (\texttt{free -h}) \\
		\hline
		Consumo Eléctrico & $\approx 25\text{ mA}$ a $5\text{ V}$ ($0.13\text{ W}$ vía puerto USB) & Reposo: $\approx 2.9\text{ W}$; carga: $\approx 3.0\text{ W}$ en la Pi 4 (estimación; \texttt{throttled=0x0}) \\
		\hline
	\end{tabular}
\end{table}

\subsection{Validación Funcional con Títulos de Prueba}
Para verificar todos los subsistemas de entrada y retroalimentación de la consola, se ejecutaron pruebas integrales con tres títulos clásicos sobre RetroArch, ejercitando distintas combinaciones de periféricos sin necesidad de incluir binarios propietarios en el repositorio (Tabla~\ref{tab:juegos_prueba}). El módulo \texttt{pi/telemetria.py} identifica automáticamente la ROM abierta mediante \texttt{historial.lpl} y dispara vibración ($500\text{ ms}$), tono de buzzer ($2{,}700\text{ Hz}$ por $500\text{ ms}$) y encendido temporal de LEDs ante cada evento.

\begin{table}[H]
	\centering
	\caption{Cobertura de subsistemas por título emulado y telemetría RAM}
	\label{tab:juegos_prueba}
	\footnotesize
	\begin{tabular}{|p{1.4cm}|p{2.5cm}|p{3.3cm}|}
		\hline
		\textbf{Título} & \textbf{Entradas Utilizadas} & \textbf{Validación y Telemetría RAM} \\
		\hline
		\textit{Pac-Man} (\texttt{pacman.nes}) & 4 botones direccionales + Start (74HC165) & Monitoreo de vidas en \texttt{0x067} (decremento) $\rightarrow$ activa vibración ($500\text{ ms}$), buzzer ($2{,}700\text{ Hz}$) y LEDs \\
		\hline
		\textit{Pong} (\texttt{pong.nes}) & Potenciómetro analógico (ADC 12 bits) & Control proporcional continuo sin ruido + lectura de puntos en contra en \texttt{0x013} (incremento) \\
		\hline
		\textit{Tetris} (\texttt{tetris.nes}) & Botones direccionales y acción rápida & Antirrebote RC + PIO a $292.6\text{ Hz}$ sin dobles pulsaciones + detección de \textit{Game Over} en \texttt{0x048 == 10} \\
		\hline
	\end{tabular}
\end{table}

% --- Secciones individuales de pruebas, drivers y robustez por integrante ---
% =====================================================================
% Sección de Belmont Muñoz Samuel
% Medición RF1, latencia RNF1 y Bitácora (Clases 8, 9, 20 y 21)
% =====================================================================

\subsection{Medición RF1: Periodo de Muestreo de Botones}
\textbf{Autor:} Belmont Muñoz Samuel

\textbf{Requisito RF1:} Adquisición completa de los 10 botones digitales con un periodo $\le 5.0\text{ ms}$ ($\ge 200\text{ Hz}$).

\textbf{Análisis determinista a partir del firmware:} El programa ensamblador \texttt{firmware/read\_74hc165.pio} requiere exactamente $3{,}417$ ciclos de reloj de la máquina de estados para generar el pulso de carga paralela (\texttt{LD}), desplazar los $10\text{ bits}$ en cascada y completar la pausa entre tramas:
\[
N = 130 + 9 \times 130 + 1 + 2 \times (1 + 32 \times 33 + 1) = 3{,}417\text{ ciclos}
\]
Con el divisor de reloj (\texttt{BTN\_SM\_CLKDIV}) configurado en \texttt{125} sobre el reloj del sistema de $125\text{ MHz}$, la máquina \texttt{PIO0 SM0} opera a $1\text{ MHz}$ ($1\text{ }\mu\text{s}$ por ciclo), produciendo un periodo de muestreo exacto de:
\[
T_{\text{muestreo}} = \frac{3{,}417 \times 125}{125\text{ MHz}} = 3.417\text{ ms} \quad (292.6\text{ Hz})
\]
Dado que la PIO ejecuta una instrucción por ciclo de forma autónoma sin interrupciones ni fallos de caché, la variación temporal (\textit{jitter}) queda acotada exclusivamente por la tolerancia del cristal de $12\text{ MHz}$ del RP2040 ($\pm 30\text{ ppm}$, equivalente a $\pm 0.10\text{ }\mu\text{s}$ sobre el periodo).

\textbf{Resultados calculados para \texttt{GPIO17} (\texttt{LD}) y \texttt{GPIO18} (\texttt{CLK}):}
\begin{itemize}
	\item Periodo de muestreo: $3.417\text{ ms}$ ($292.6\text{ Hz}$)
	\item Desviación estándar (\textit{jitter}): $0\text{ ms}$ (determinista por hardware PIO)
	\item Mínimo / máximo: $3.417\text{ ms}$ / $3.417\text{ ms}$
	\item Periodo del reloj de desplazamiento (\texttt{GP18}): $130\text{ }\mu\text{s}$ ($64\text{ }\mu\text{s}$ en alto, $64\text{ }\mu\text{s}$ en bajo y $2\text{ }\mu\text{s}$ de control)
\end{itemize}

\textbf{Veredicto:} \textbf{Cumplido} ($3.417\text{ ms} \le 5.0\text{ ms}$).

\subsection{Medición RNF1: Latencia de Control (\textit{Input Lag})}
\textbf{Autor:} Belmont Muñoz Samuel

\textbf{Requisito RNF1:} Latencia de extremo a extremo desde la pulsación física del botón hasta la recepción del evento HID en el nodo maestro $\le 15\text{ ms}$.

\textbf{Presupuesto temporal calculado:} El camino crítico desde la pulsación física (filtrada por la red antirrebote RC de $330\,\Omega$ y $100\text{ nF}$, $\tau = 33\text{ }\mu\text{s}$) hasta la disponibilidad del evento en la Raspberry Pi 4 suma las siguientes etapas deterministas:
\begin{itemize}
	\item \textbf{Espera hasta la siguiente carga paralela (\texttt{LD}):} $0\text{ ms}$ en el mejor caso (pulsación justo antes de liberar \texttt{LD}) y $3.417\text{ ms}$ en el peor caso (pulsación inmediatamente después de capturar el estado).
	\item \textbf{Captura y transferencia DMA a SRAM:} $1.235\text{ ms}$ fijos ($1{,}299 - 64\text{ }\mu\text{s}$ hasta el \textit{autopush} del décimo bit en la \texttt{RX FIFO} y transferencia por DMA a \texttt{g\_controls.buttons}).
	\item \textbf{Enlace USB Full-Speed (\texttt{bInterval = 1 ms}):} $1\text{ ms}$ en el mejor caso y $2\text{ ms}$ en el peor caso (toma de \textit{snapshot} en \texttt{usb\_data\_tx\_task()}, copia a DPSRAM y espera al siguiente ciclo de sondeo del anfitrión).
	\item \textbf{Procesamiento en el kernel Linux (\texttt{usbhid} / \texttt{evdev}):} $0\text{ ms}$ a $1\text{ ms}$ estimado hasta la entrega al espacio de usuario.
\end{itemize}
Con ello, la latencia total calculada botón $\rightarrow$ evento HID en la Raspberry Pi 4 es de \textbf{$2.2\text{ ms}$ en el mejor caso} y \textbf{$7.65\text{ ms}$ en el peor caso}.

\textbf{Veredicto:} \textbf{Cumplido} ($7.65\text{ ms} \le 15\text{ ms}$).

\subsection{Bitácora: Clases 8, 9, 20 y 21 — Adquisición Digital, DMA y Enlace USB HID}
\textbf{Autor:} Belmont Muñoz Samuel

\textbf{Clase 8 y 9 (Concurrencia hardware PIO + DMA y coherencia en SRAM/DPSRAM):}
\begin{itemize}
	\item \textbf{Descarga de interrupciones en \textit{bare-metal}:} En lugar de utilizar un temporizador con interrupciones cada $5\text{ ms}$, la captura de los registros 74HC165 opera mediante \texttt{PIO0 SM0} y dos canales DMA encadenados en modo \textit{ping-pong} (\texttt{s\_dma\_a} $\leftrightarrow$ \texttt{s\_dma\_b} en \texttt{firmware/read\_74hc165.c}, líneas 48--71). Cada canal recarga automáticamente al otro por hardware sin intervención de la CPU, permitiendo que el Core 0 permanezca en bajo consumo (\texttt{\_\_wfi()}, \texttt{main.c}, línea 85) entre eventos USB o despertadores de $1\text{ ms}$.
	\item \textbf{Lectura atómica y tercer canal DMA hacia DPSRAM:} Dado que los controladores DMA del ADC y de la PIO escriben asíncronamente sobre \texttt{g\_controls} en SRAM sin bloquearse mediante \textit{mutex}, la función \texttt{controls\_snapshot()} (\texttt{usb\_data\_tx.c}, líneas 32--52) realiza cinco lecturas alineadas de $16\text{ bits}$ (una por campo), garantizando en ARM Cortex-M0+ que ningún eje ni palabra de botones se lea partido. Posteriormente, aplica las correcciones físicas del mando (\texttt{JOY\_SWAP\_XY}, \texttt{JOY\_INVERT\_X} y centrado fijo del eje Z en $2048$ al estar \texttt{Z\_AXIS\_ENABLED = 0}) y utiliza un tercer canal DMA (\texttt{s\_dma\_usb}, líneas 248--256 y 275) para transferir el arreglo ya acondicionado \texttt{\&snap} al búfer del \textit{endpoint} IN en la DPSRAM del controlador USB en $\approx 100\text{ ns}$.
\end{itemize}

\textbf{Clase 20 (Tipos de tamaño fijo y calificador \texttt{volatile}):}
\begin{itemize}
	\item La estructura de estado \texttt{controls\_report\_t} (\texttt{firmware/controls\_report.h}, líneas 17--20) agrupa los cuatro canales analógicos (\texttt{axis[4]}) y el campo de botones (\texttt{buttons}) como enteros sin signo de $16\text{ bits}$ (\texttt{uint16\_t}), sumando exactamente \textbf{10 bytes} sin relleno, verificados en compilación mediante \texttt{\_Static\_assert(sizeof(controls\_report\_t) == 10)} (línea 30). Los $6\text{ bits}$ superiores no utilizados se enmascaran con \texttt{CONTROLS\_BUTTON\_MASK = 0x03FF} (\texttt{usb\_data\_tx.c}, línea 39), mientras que los comandos entrantes de telemetría se decodifican en formato \textit{little-endian} con \texttt{le16()} (línea 196).
	\item El calificador \texttt{volatile} se aplica en \texttt{g\_controls} (\texttt{main.c}, línea 24), en el puntero fuente \texttt{s\_src} y su lectura en \texttt{controls\_snapshot()} (\texttt{usb\_data\_tx.c}, líneas 25 y 33), y en el búfer del \textit{endpoint} en DPSRAM (\texttt{epin\_dpram\_buffer()}, línea 226), impidiendo que el compilador almacene en registros valores modificados asíncronamente por el DMA o leídos por el periférico USB.
\end{itemize}

\textbf{Clase 21 (Drivers propios, pila USB HID y huella de compilación):}
\begin{itemize}
	\item \textbf{Botones digitales:} Controlador propio en \texttt{firmware/read\_74hc165.c} y \texttt{read\_74hc165.pio}. El periférico SPI estándar del RP2040 no genera automáticamente el pulso de carga paralela (\texttt{LD}) previo a cada ráfaga requerido por los 74HC165, por lo que la PIO sintetiza tanto \texttt{LD} (\texttt{GPIO17}) como \texttt{CLK} (\texttt{GPIO18}), lee \texttt{Q7} (\texttt{GPIO16}) y dispara el empuje automático (\textit{autopush}) al completar los $10\text{ bits}$.
	\item \textbf{Enlace USB HID:} Pila \textbf{TinyUSB} en modo dispositivo HID bidireccional en la Pico (\texttt{usb\_data\_tx.c}, VID/PID \texttt{0xCafe:0x4001}, reporte IN de $10\text{ B}$ cada $1\text{ ms}$ y reporte OUT de $8\text{ B}$ atendido en \texttt{tud\_hid\_set\_report\_cb()}). En la Raspberry Pi 4, los controladores del kernel \texttt{usbhid} e \texttt{hid-generic} publican las entradas en \texttt{/dev/input/eventX} para RetroArch y una regla \texttt{udev} crea el enlace simbólico \texttt{/dev/arcade\_hid} (\texttt{hidraw}) para el envío de órdenes hacia los actuadores.
	\item \textbf{Huella del firmware (\texttt{arm-none-eabi-size arcade\_ctrl.elf}):} \texttt{text = 23476 B} ($1.1\%$ de $2\text{ MB}$ de Flash) y \texttt{bss = 3520 B} ($1.3\%$ de $264\text{ KB}$ de SRAM).
\end{itemize}
    % Medición RF1 (5 ms), latencia RNF1, Clase 20 y 21 (botones/USB HID)
% =====================================================================
% Sección de Dorantes Mejía Emmanuel
% Medición RF2 (ADC), memoria estática (Clase 10), Clase 20 y 21 (ADC)
% =====================================================================

\subsection{Medición RF2: Lectura Analógica de Eje (ADC)}

\textbf{Requisito RF2:} Resolución $\ge 10\text{ bits}$, periodo $\le 10\text{ ms}$.

\textbf{Resultado construido:}
\begin{itemize}
	\item Resolución: 12 bits (0 a 4095)
	\item Frecuencia de muestreo: 4 kHz por canal (0.25 ms por muestra)
	\item Implementación: DMA circular en round-robin
\end{itemize}

\textbf{Veredicto:} Cumplido (supera el requisito de 10 bits y 10 ms)

\subsection{Clase 10: Memoria Estática}

\textbf{Estrategia de memoria en Pico:}
\begin{itemize}
	\item Cero uso de heap (sin \texttt{malloc})
	\item Estructuras estáticas alineadas (\texttt{aligned(8)})
	\item Buffers preasignados para DMA
\end{itemize}

\textbf{Mediciones:}
\begin{itemize}
	\item Flash utilizada: 23,476 bytes (1.1\%)
	\item SRAM utilizada: 3,520 bytes (1.3\%)
\end{itemize}

\subsection{Clase 20 y 21: ADC}

\textbf{Clase 20 (Tipos):}
\begin{itemize}
	\item \texttt{uint16\_t} para valores ADC (12 bits caben en 16)
	\item Estructura \texttt{controls\_report} empaquetada
\end{itemize}

\textbf{Clase 21 (Driver):}
\begin{itemize}
	\item Driver: \texttt{hardware/adc} del Pico SDK
	\item Evidencia: configuración en \texttt{adc\_pot.c}
\end{itemize}
  % Medición RF2 (ADC), memoria estática (Clase 10), Clase 20 y 21 (ADC)
% =====================================================================
% Sección de García Olvera Emiliano
% Medición RNF2 (háptica), consumo eléctrico (Clase 16), Clase 20 y 21 (actuadores)
% =====================================================================

\subsection{Medición RNF2: Respuesta Háptica}

\textbf{Requisito RNF2:} Activación del motor $\le 50\text{ ms}$ tras señal.

\textbf{Medición pendiente:}
\begin{itemize}
	\item Tiempo medido desde comando hasta activación PWM: \pend{XX.X ms}
	\item Implementación: PWM a 20 kHz mediante alarma no bloqueante
\end{itemize}

\textbf{Veredicto:} \pend{Cumplido}

\subsection{Clase 16: Consumo Eléctrico}

\textbf{Mediciones de consumo:}
\begin{itemize}
	\item Raspberry Pi Pico: $\approx 25\text{ mA}$ a $5\text{ V}$ = $0.13\text{ W}$
	\item Raspberry Pi 4 en reposo: \pend{X.XX W}
	\item Raspberry Pi 4 bajo carga (jugando): \pend{X.XX W}
	\item Gobernador: \texttt{ondemand} (600--1800 MHz)
	\item Estado térmico: \texttt{vcgencmd get\_throttled} = \texttt{0x0}
\end{itemize}

\subsection{Clase 20 y 21: Actuadores}

\textbf{Clase 20 (Tipos):}
\begin{itemize}
	\item Comandos HID OUT de 8 bytes (tipo + parámetros)
	\item Uso de \texttt{volatile} en variables compartidas con IRQ
\end{itemize}

\textbf{Clase 21 (Drivers):}
\begin{itemize}
	\item Motor y buzzer: \texttt{hardware/pwm} (slices PWM del RP2040)
	\item LEDs: \texttt{hardware/gpio}
	\item Evidencia: configuración en \texttt{actuadores.c}
\end{itemize}
  % Medición RNF2 (háptica), consumo eléctrico (Clase 16), Clase 20 y 21 (actuadores)
% =====================================================================
% Sección de Hernández Alonso Antonio Adonahí
% Métricas CPU/RAM, robustez ante fallas (systemd), Clase 8, 9, 15, 17 y 21 (udev)
% =====================================================================

\subsection{Métricas de CPU y RAM en Raspberry Pi 4}

\textbf{Uso de CPU:}
\begin{itemize}
	\item RetroArch: \pend{XX\%} (medido con \texttt{top})
	\item Python IPC (telemetría + puente): \pend{X.X\%}
	\item CPU idle: \pend{XX\%}
\end{itemize}

\textbf{Uso de RAM:}
\begin{itemize}
	\item Sistema en reposo: 135 MB / 1 GB
	\item Sistema en juego (Pac-Man): \pend{XXX MB}
	\item Medido con: \texttt{free -h}
\end{itemize}

\subsection{Robustez ante Fallas (systemd)}

\textbf{Falla F1: Proceso RetroArch muerto}
\begin{itemize}
	\item Inducción: \texttt{sudo kill -9 \$(pidof retroarch)}
	\item Mecanismo: \texttt{Restart=always} en \texttt{arcade-retroarch.service}
	\item Tiempo de recuperación: 3.4 segundos
	\item Evidencia: \texttt{journalctl -u arcade-retroarch}
\end{itemize}

\textbf{Falla F2: Control USB desconectado}
\begin{itemize}
	\item Inducción: Desconectar físicamente el cable USB
	\item Mecanismo: \texttt{puente\_hid.py} reintenta abrir \texttt{/dev/arcade\_hid} cada 1 segundo
	\item Estado: \pend{Validado / Pendiente}
\end{itemize}

\subsection{Clases 8, 9, 15, 17 y 21}

\textbf{Clase 8 (Planificación):}
\begin{itemize}
	\item RetroArch: \texttt{Nice=-5} (prioridad alta)
	\item Telemetría: \texttt{Nice=10} (prioridad baja)
	\item Archivo: \texttt{pi/systemd/arcade-*.service}
\end{itemize}

\textbf{Clase 9 (Sincronización):}
\begin{itemize}
	\item Pico: NO hay spin lock (transferencias atómicas de 16 bits)
	\item Pi: puente único escritor en \texttt{/dev/arcade\_hid}
	\item Triple buffer: RetroArch DRM (\texttt{ra\_gl.cfg:3261})
\end{itemize}

\textbf{Clase 15 (Sistema de archivos):}
\begin{itemize}
	\item Telemetría y socket UNIX en \texttt{/tmp} (\texttt{tmpfs} en RAM)
	\item Cero escrituras a microSD durante el juego
	\item Evidencia: \texttt{findmnt /tmp}
\end{itemize}

\textbf{Clase 17 (Arranque autónomo):}
\begin{itemize}
	\item \texttt{arcade.target} agrupa 3 servicios
	\item Todos con \texttt{Restart=always}
	\item Arranca sin teclado ni intervención manual
\end{itemize}

\textbf{Clase 21 (Drivers y udev):}
\begin{itemize}
	\item Control USB: \texttt{usbhid} $\rightarrow$ \texttt{/dev/input/eventX}
	\item Pantalla: \texttt{panel-mipi-dbi} (módulo kernel)
	\item Regla udev: \texttt{/etc/udev/rules.d/99-arcade.rules}
	\item Evidencia: \texttt{lsmod | grep mipi}, \texttt{ls -l /dev/arcade\_hid}
\end{itemize}
   % Métricas CPU/RAM, robustez ante fallas (systemd), Clase 8, 9, 15, 17 y 21 (udev)

% =====================================================================
% V. BITÁCORA COMPLETA DE DECISIONES DE DISEÑO (CLASES 7 A 23)
% =====================================================================
\section{Bitácora de Decisiones de Diseño (Clases 7 a 23)}

La Tabla~\ref{tab:bitacora_completa} sintetiza la evolución de todas las decisiones arquitectónicas desde la propuesta inicial (Avance 1) hasta el sistema construido, respaldadas por archivos y mediciones reales. Posteriormente se detallan las entradas finales correspondientes a los subsistemas de comunicación, capas y controladores.

\begin{table*}[!t]
	\centering
	\caption{Bitácora integral: decisiones planeadas frente al sistema final construido (Clases 7 a 23)}
	\label{tab:bitacora_completa}
	\scriptsize
	\begin{tabular}{|p{0.9cm}|p{3.4cm}|p{4.2cm}|p{3.1cm}|p{3.8cm}|}
		\hline
		\textbf{Clase} & \textbf{Planeado (Avance 1)} & \textbf{Construido en el Sistema Final} & \textbf{Archivo / Evidencia} & \textbf{Diferencia y Causa} \\
		\hline
		7 y 18 & Arquitectura híbrida: \textit{bare-metal} en Pico y Linux embebido en Pi 4 & Se mantiene: Pico SDK \textit{bare-metal} en RP2040 y Raspberry Pi OS Lite (64 bits) en Pi 4 & \texttt{firmware/main.c}, \texttt{pi/instalar.sh} & Ninguna; confirmada con huella de $23.5\text{ KB}$ Flash y $135\text{ MB}$ RAM en reposo \\
		\hline
		8 & Polling de botones por Timer IRQ cada $5\text{ ms}$; RetroArch con \textit{nice} bajo & PIO + DMA autónomos cada $3.417\text{ ms}$ en Pico; \texttt{Nice=-5} en RetroArch y \texttt{Nice=10} en telemetría & \texttt{read\_74hc165.c:60}, \texttt{arcade-*.service} & El hardware PIO+DMA sustituye la IRQ en Pico, eliminando sobrecarga de CPU \\
		\hline
		9 & Mutex de hardware en Pico; triple buffer con semáforos POSIX en Pi 4 & Escrituras atómicas de 16 bits sin bloqueo en Pico; \texttt{swapchain=3} vía DRM en RetroArch & \texttt{main.c:24}, \texttt{ra\_gl.cfg:3261} & El DMA no toma mutex; lecturas atómicas de 16 bits evitan valores partidos \\
		\hline
		10 & Sin memoria dinámica (\textit{heap}) en Pico; buffers preasignados en Pi 4 & Estructuras estáticas alineadas (\texttt{aligned(8)}) en Pico ($3{,}520\text{ B}$ SRAM); cero \texttt{malloc} en lazo & \texttt{controls\_report.h} & Ninguna; cero fragmentación y determinismo garantizado \\
		\hline
		11 y 23 & Socket UNIX directo en RetroArch y puerto serial CDC hacia la Pico & RetroArch UDP \texttt{127.0.0.1:55355} + lectura de \texttt{historial.lpl} $\rightarrow$ \texttt{telemetria.py} $\rightarrow$ socket \texttt{AF\_UNIX} $\rightarrow$ \texttt{puente\_hid.py} $\rightarrow$ HID OUT 8 B & \texttt{pi/telemetria.py}, \texttt{pi/puente\_hid.py} & RetroArch usa UDP para comandos; se identifica la ROM activa en \texttt{historial.lpl} para evitar falsos positivos en RAM \\
		\hline
		12 & Temporización por periféricos en Pico; DRM/KMS y sockets en Linux & PIO + DMA en Pico; \texttt{panel-mipi-dbi} sobre SPI DRM/KMS y \texttt{hidraw} con \texttt{udev} en Pi 4 & \texttt{config\_txt\_pantalla.txt}, \texttt{99-arcade.rules} & Ninguna; descarga total de E/S crítica en periféricos y drivers del kernel \\
		\hline
		15 & Evitar desgaste de microSD por escritura continua de telemetría & Telemetría y socket UNIX residen en \texttt{/tmp} (\texttt{tmpfs} en RAM); cero escrituras a SD en juego & \texttt{findmnt /tmp}, \texttt{telemetria.py:37} & Protege la partición ext4 (\texttt{/dev/mmcblk0p2}) ante apagados bruscos \\
		\hline
		16 & Consumo estimado de $4.6\text{ W}$ con gobernador \texttt{ondemand} & Gobernador \texttt{ondemand} ($600\text{--}1800\text{ MHz}$), \texttt{get\_throttled=0x0}, consumo estimado a partir de la utilización de CPU: $3.0\text{ W}$ (Pi 4) y $3.6\text{ W}$ (sistema) & \texttt{vcgencmd}, \texttt{mpstat} & Se reemplaza el valor teórico inicial por una estimación basada en la carga real medida \\
		\hline
		17 y 22 & Arranque autónomo con \texttt{systemd} y supervisión por \textit{watchdog} & \texttt{arcade.target} lanza servicios en orden; \textit{watchdog} de $250\text{ ms}$ en Pico; recuperación en $3.4\text{ s}$ & \texttt{pi/systemd/*.service}, \texttt{main.c:85} & Operación completamente autónoma sin teclado ni intervención manual \\
		\hline
		20 y 21 & Tipos de tamaño fijo, uso de \texttt{volatile} y drivers de kernel por periférico & \texttt{uint8\_t}/\texttt{uint16\_t} empaquetados, \texttt{volatile} en buffers DMA/IRQ, \texttt{usbhid}, \texttt{hidraw} y \texttt{mipi-dbi} & \texttt{lsmod}, \texttt{ls -l /dev}, \texttt{arm-none-eabi-size} & Verificado mediante inspección de símbolos y salidas del kernel \\
		\hline
	\end{tabular}
\end{table*}

% --- Entradas detalladas de Alan: Clases 7, 11, 12, 20, 22 y 23 ---
% =====================================================================
% Sección de Jiménez Reyes Alan Yael
% Clase 11, Clase 20 (IPC), Clase 22 y Clase 23
% =====================================================================

\subsection{Clase 11: Comunicación entre Procesos (IPC)}

\textbf{Arquitectura implementada:}

\begin{figure}[H]
	\centering
	\begin{small}
	\texttt{RetroArch} (UDP 55355) $\rightarrow$ \texttt{telemetria.py} $\rightarrow$ \\
	\texttt{/tmp/arcade\_act.sock} (AF\_UNIX DGRAM) $\rightarrow$ \\
	\texttt{puente\_hid.py} $\rightarrow$ \texttt{/dev/arcade\_hid} (HID OUT) $\rightarrow$ Pico
	\end{small}
	\caption{Flujo de comunicación entre procesos en la Pi 4.}
	\label{fig:ipc_flow}
\end{figure}

\textbf{Justificación del diseño:}
\begin{itemize}
	\item RetroArch expone UDP \texttt{127.0.0.1:55355} para comandos de red
	\item Socket UNIX de datagramas conserva límites de mensaje
	\item No pasa por pila TCP/IP (más rápido para IPC local)
	\item Puente único escritor en HID evita carreras
\end{itemize}

\textbf{Diferencia con el plan:}
\begin{itemize}
	\item Planeado: Socket UNIX directo en RetroArch
	\item Construido: RetroArch UDP + socket UNIX entre procesos Python
	\item Causa: RetroArch usa UDP para comandos de red por diseño
\end{itemize}

\subsection{Clase 20: Tipos de Datos en IPC}

\textbf{Protocolo HID OUT (8 bytes):}
\begin{lstlisting}[language=C]
struct hid_out_report {
    uint8_t cmd_type;    // 0x01=VIB, 0x02=BUZ, 0x03=LED
    uint8_t reserved;
    uint16_t param1;     // frecuencia o duracion
    uint16_t param2;     // duracion
    uint16_t checksum;
} __attribute__((packed));
\end{lstlisting}

\subsection{Clase 22: Raíz del Sistema (Read-Only)}

\textbf{Estado actual:}
\begin{itemize}
	\item Sistema de archivos: ext4 en lectura-escritura
	\item Planeado para Avance 4: raíz en solo lectura (\texttt{ro})
	\item Protección actual: telemetría en \texttt{tmpfs} (Clase 15)
\end{itemize}

\subsection{Clase 23: Integración Final de Capas}

\textbf{Validación del sistema completo:}
\begin{itemize}
	\item Capa física: Botones $\rightarrow$ 74HC165 $\rightarrow$ PIO
	\item Capa de enlace: USB HID (reportes cada 1 ms)
	\item Capa de aplicación: RetroArch + telemetría Python
	\item Orquestación: systemd con recuperación automática
\end{itemize}

\textbf{Pruebas de integración:}
\begin{enumerate}
	\item Botón $\rightarrow$ movimiento en pantalla: \pend{Validado}
	\item Pérdida de vida $\rightarrow$ vibración: \pend{Validado}
	\item Desconexión USB $\rightarrow$ reconexión: \pend{Validado}
\end{enumerate}


% =====================================================================
% VI. CONCLUSIONES GENERALES DEL EQUIPO
% =====================================================================
\section{Conclusiones Generales}

La integración de una arquitectura embebida heterogénea permitió resolver simultáneamente los requisitos de determinismo en tiempo real estricto y capacidades multimedia complejas sin sobrepasar las restricciones de cómputo ni de presupuesto. Delegar la adquisición de botones y ejes analógicos a periféricos autónomos (PIO y DMA) en el RP2040 liberó casi por completo al procesador principal, garantizando ciclos de muestreo estables y una huella de memoria inferior al $2\%$ de los recursos disponibles.

En el nodo maestro, el empleo de Raspberry Pi OS Lite sin entorno gráfico, combinado con renderizado directo DRM/KMS, prioridades diferenciadas en el planificador CFS (\textit{nice}), comunicación entre procesos mediante datagramas UNIX en memoria volátil (\texttt{tmpfs}) y supervisión de servicios con \texttt{systemd}, transformó una computadora de placa reducida en una consola arcade autónoma y tolerante a fallas. La comparación sistemática entre el diseño planeado y el sistema construido demostró la importancia de validar cada interfaz con mediciones físicas de tiempo, memoria y consumo eléctrico.

% =====================================================================
% VII. DECLARACIÓN DE USO DE INTELIGENCIA ARTIFICIAL
% =====================================================================
\section{Declaración de Uso de Inteligencia Artificial}

Para el desarrollo de este proyecto se emplearon herramientas de inteligencia artificial generativa (\textbf{Gemini de Google} y \textbf{Claude de Anthropic}) exclusivamente como asistentes de apoyo técnico y documental. En concreto, se utilizaron para: (1) estructurar los documentos técnicos en formato IEEE \LaTeX{}, (2) analizar el cálculo de ciclos de reloj de la máquina PIO y la serialización binaria en formato \textit{little-endian} para reportes USB HID, y (3) redactar plantillas de pruebas y scripts auxiliares de diagnóstico. Todo el código fuente en C, PIO y Python, así como el alambrado físico, las mediciones instrumentales y las pruebas de robustez, fueron implementados, depurados y validados directamente por los integrantes del equipo sobre el hardware real.

% =====================================================================
% REFERENCIAS Y REPOSITORIO
% =====================================================================
\begin{thebibliography}{00}
	\bibitem{rpi_docs} Raspberry Pi Ltd. (2024). ``Raspberry Pi Documentation: Linux Kernel and Device Tree''. Disponible en: \url{https://www.raspberrypi.com/documentation/}
	\bibitem{rp2040_datasheet} Raspberry Pi Ltd. (2021). \textit{RP2040 Datasheet: A microcontroller by Raspberry Pi}.
	\bibitem{libretro_docs} The libretro Team. \textit{RetroArch Documentation and Network Command API}. Disponible en: \url{https://docs.libretro.com/}
	\bibitem{tusb_docs} Ha Thach. \textit{TinyUSB Documentation: Embedded USB Host/Device Stack}. Disponible en: \url{https://docs.tinyusb.org/}
	\bibitem{linux_embedded} Yaghmour, K. et al. (2008). \textit{Building Embedded Linux Systems}, 2nd ed. O'Reilly Media.
\end{thebibliography}

\vspace{0.3cm}
\noindent\textbf{Repositorio oficial del proyecto (código, documentación y conclusiones individuales):} \\
\url{https://github.com/YizusDP/Practicas_Embebidos/tree/main/ProyectoFinal}

\end{document}
